% !TEX root = main.tex
\documentclass[main.tex]{subfiles}
\begin{document}

\section{Constrained Power Diagrams}
\label{sec:contributions}

In this section, we present the construction of power diagrams whose cells have prescribed centroids $\smash{\tcentroid_i \!\in\! \R^d}$, collected as $\smash{\tcentroid \!=\! \{\tcentroid_i\}_{i=1}^n}$. %, chosen independently of the sites.
Similar to the volume-constrained power diagrams in \eqref{eq:volume-dual}, we show that these centroid-constrained power diagrams are generated by maximizing an unconstrained concave objective, which yields a unique partition whenever one exists, both for centroid constraints alone (\S\ref{sec:centroid-constraints}) and jointly with volume constraints (\S\ref{sec:joint-constraints}).
We further propose a softening of the centroid constraints to handle targets that no power diagram realizes (\S\ref{sec:regularized-joint}), extend this softening to volume constraints (\S\ref{sec:softened-volume}), and relate our results to prior work (\S\ref{sec:remarks}).

\subsection{Centroid Constraints}
\label{sec:centroid-constraints}

We first consider the construction of a partition $\cell$ associated with fixed weighted sites $(\qq, \weight)$ whose cells have prescribed centroids $\tcentroid$.
Reusing the objective from \eqref{eq:power-diagram}, we formulate this problem as
\begin{subequations}
\label{eq:centroid-constrained}
\begin{alignat}{2}
  & \min_{\cell} \quad && \sum\nolimits_{i} \int_{\cell_i} \bigl( \pdist{\xx}{i} \bigr) \, d\xx \label{eq:centroid-constrained-obj} \\
  & \text{s.t.} \quad && \int_{\cell_i} (\xx - \tcentroid_i) \, d\xx = 0 \ \ \forall i. \label{eq:centroid-constrained-cons}
\end{alignat}
\end{subequations}
% 
Introducing a vector multiplier $\shift_i \!\in\! \R^d$, which we call the \emph{shift}, for each centroid constraint \eqref{eq:centroid-constrained-cons} yields the Lagrangian
\begin{equation}
\lag(\cell, \shift)
\!=\! 
\sum\nolimits_{i} \int_{\cell_i}\! \bigl( \pdist{\xx}{i} \bigr) d\xx - 2 \shift_i\tr\! \int_{\cell_i}\! (\xx - \tcentroid_i) \, d\xx,
\label{eq:centroid-lagrangian}
\end{equation}
% 
where we indicate the collected shifts as $\smash{\shift\!=\!\{\shift_i\}_{i=1}^n}$ and scale each $\shift_i$ by a factor of two to make our derivatives more concise.
% 
For fixed shifts, the combined integrand in \eqref{eq:centroid-lagrangian} can be rearranged as
\begin{equation}
  \pdist{\xx}{i} - 2\,\shift_i\tr (\xx - \tcentroid_i)
    = \norm{\xx - \dual_i}^2 - \combweight_i ,
  \label{eq:complete-square}
\end{equation}
based on the \emph{shifted} weighted sites $(\dual_i, \combweight_i)$ of the form
\begin{subequations}
  \label{eq:shifted-site}
  \begin{empheq}[left={\empheqlbrace\,}]{align}
    \dual_i &= \qq_i + \shift_i , \label{eq:shifted-site-pos} 
    \\[2mm]
    \combweight_i &= \weight_i + \norm{\shift_i}^2 + 2\,\shift_i\tr (\qq_i - \tcentroid_i). \label{eq:shifted-site-weight}
  \end{empheq}
\end{subequations}
% 
We thus conclude that, for fixed shifts, minimizing \eqref{eq:centroid-lagrangian} over partitions yields the power diagram of the shifted weighted sites \eqref{eq:shifted-site}.
% 
We call the shifted sites $\dual$ the \emph{generators} of this power diagram, to distinguish them from the sites $\qq$ used in the objective \eqref{eq:centroid-constrained-obj}.
Importantly, the properties of power diagrams reviewed in \S\ref{sec:background}, such as the homotheties $\homotheties$ and the normals of $\achievable$, apply to the generators $\dual$ rather than to the sites $\qq$.
% 
As in \S\ref{sec:background}, we restrict the shifts to configurations with pairwise distinct generators $\dual_i$, which guarantees that this minimizer is a unique power diagram. % on which $\centroiddual$ is differentiable.

The optimal value of \eqref{eq:centroid-lagrangian} defines the \emph{dual objective}
\begin{equation}
  \centroiddual(\shift)
  % =
  % \centroiddual(\shift; \qq, \weight, \tcentroid)
  = 
  \min_{\cell} \lag(\cell, \shift)
  = 
  \energy\bigl( \dual(\shift), \combweight(\shift) \bigr).
  \label{eq:centroid-dual}
\end{equation}
The constrained optimization in \eqref{eq:centroid-constrained} thus becomes the {unconstrained} maximization of $\centroiddual$ over the space of power diagrams, parametrized by the shifts $\shift$.
% , with the weighted sites $(\qq, \weight)$ and the targets $\tcentroid$ as fixed inputs.
% 
As detailed in Appendix~\ref{sec:appendix-centroid}, the gradient of $\centroiddual$ recovers the centroid constraint
\begin{equation}
  \partial_{\shift_i} \centroiddual
    = 2 \int_{\cell_i} (\tcentroid_i - \xx) \, d\xx
    = 2 \, \cellvol_i \, (\tcentroid_i - \centroid_i) ,
  \label{eq:centroid-dual-gradient}
\end{equation}
so a stationary shift is exactly one whose cells realize the prescribed centroids.
The Hessian of $\centroiddual$ is a facet-based vector Laplacian coupling neighboring shifts through the facets $\face_{ij}$ of the power diagram with $d \!\times\! d$ blocks
\begin{subequations}
  \label{eq:centroid-hessian}
  \begin{empheq}[left={\empheqlbrace\,}]{align}
    \partial^2_{\shift_i \shift_j}
    \centroiddual &= \frac{2}{\norm{\dual_i - \dual_j}} \int_{\face_{ij}} (\xx - \tcentroid_i)(\xx - \tcentroid_j)\tr \, dA,
    \label{eq:centroid-hessian-off}
    \\[2mm]
    \partial^2_{\shift_i \shift_i}
    \centroiddual &= -\!\sum\nolimits_{j \sim i} \frac{2}{\norm{\dual_i - \dual_j}} \int_{\face_{ij}} (\xx - \tcentroid_i)(\xx - \tcentroid_i)\tr \, dA.
    \label{eq:centroid-hessian-diag}
  \end{empheq}
\end{subequations}
In Appendix~\ref{sec:appendix-centroid}, we also show that \eqref{eq:centroid-hessian} is \emph{negative semi-definite} and, hence, $\centroiddual$ is concave with respect to shifts.
Therefore, the centroid-constrained problem is an \emph{unconstrained concave maximization}.

It is worth noting that the Hessian null-space is not always trivial and it may contain infinitesimal homotheties of the sites that the shifts can realize while keeping the weights fixed.
To express these null-space modes explicitly, we associate each site with the scalar
\begin{equation}
  \gamma_i = \norm{\tcentroid_i}^2 - \bigl( \norm{\tcentroid_i - \qq_i}^2 - \weight_i \bigr),
  \label{eq:pivot}
\end{equation}
which is the height at the target $\tcentroid_i$ for the hyperplane lifted from the weighted site $(\qq_i, \weight_i)$.
As shown in Appendix~\ref{sec:appendix-centroid}, the null-space consists exactly of the perturbations $\smash{\delta\shift_i = \mu\,\dual_i + \translation}$ for which
\begin{equation}
  \mu\,\gamma_i + 2\,\translation\tr\tcentroid_i + c = 0 \ \ \forall i
  \label{eq:centroid-kernel}
\end{equation}
holds for some $\smash{(\mu, \translation, c) \in \R \times \R^d \times \R}$, which are the infinitesimal dilation, translation and vertical offset of the homotheties $\homotheties$ in \eqref{eq:lifted-homothety}, with $\smash{s \!=\! 1 \!+\! \mu}$. %, which dilate and translate the shifted sites.
Geometrically, each shift $\shift_i$ tilts the lifted hyperplane of its site about the fixed {pivot} $\smash{(\tcentroid_i, \gamma_i)}$. %, since the value of this hyperplane at the target $\tcentroid_i$ does not depend on $\shift_i$.
Therefore, a perturbation $\delta\shift$ leaves the diagram unchanged only if the induced rotations amount to a common homothety of all lifted hyperplanes, which, by \eqref{eq:centroid-kernel}, requires all pivots to lie on a common hyperplane.
{Hence, $\centroiddual$ is strictly concave if and only if its pivots are not coplanar.}
% Since $n$ points in $\R^{d+1}$ always lie on a common hyperplane when $n \!\le\! d\!+\!1$, this requires $n \!>\! d\!+\!1$.
% The pivots are also coplanar in two notable configurations of the targets: when the targets lie on a common hyperplane of $\R^d$, since the pivots then lie on a vertical hyperplane over it, and, in the special setting with $\qq_i \!=\! \tcentroid_i$ and $\weight_i \!=\! 0$, when the targets lie on a common sphere, since the pivots $\smash{(\tcentroid_i, \norm{\tcentroid_i}^2)}$ then lie on the paraboloid, which a hyperplane cuts along a sphere.
Otherwise, the maximizing shifts $\shift$, if they exist, form an affine family along the modes of \eqref{eq:centroid-kernel}.
% 
Importantly, observe that this condition concerns the uniqueness of the \emph{shifts}, not of the \emph{cells}.
In both cases, the resulting power diagram is unique if it exists.
% any partition solving \eqref{eq:centroid-constrained} minimizes the Lagrangian \eqref{eq:centroid-lagrangian} at every maximizing shift, and this minimizer is unique.
% We discuss the existence of maximizers, and of maximizers with distinct shifted sites, in \S\ref{sec:remarks}.

\subsection{Volume + Centroid Constraints}
\label{sec:joint-constraints}

Next, we consider the partition that enforces both prescribed volumes $\smash{\pvol}$ \emph{and} prescribed centroids $\smash{\tcentroid}$ simultaneously.
Repeating the steps of \S\ref{sec:centroid-constraints}, we introduce a shift $\shift_i \!\in\! \R^d$ for each centroid constraint and a weight $\weight_i \!\in\! \R$ for each volume constraint, which together yield the Lagrangian
\begin{equation}
  \lag(\cell, \shift, \weight)
  \!=\! 
  \sum\nolimits_{i} 
  \!
  \int_{\cell_i} \!\!\!\!\norm{\xx -\! \qq_i}^2 \!d\xx 
  - 2\shift_i\tr 
  \!\!
  \int_{\cell_i}\!\!\!\! (\xx - \tcentroid_i) d\xx 
  - \weight_i \bigl( \cellvol_i - \pvol_i \bigr).
  \label{eq:joint-lagrangian}
\end{equation}
For fixed multipliers, the partition minimizer corresponds once again to a power diagram with shifted weighted sites $(\dual, \combweight)$ given in \eqref{eq:shifted-site}, under the same restriction to distinct shifted sites with non-empty cells as in \S\ref{sec:centroid-constraints}.
Its optimal value defines the \emph{dual objective}
\begin{equation}
  \jointdual(\shift, \weight)
  % \jointdual(\shift, \weight; \qq, \tcentroid)
    \!=\! 
    \min_{\cell} \lag(\cell, \shift, \weight)
    \!=\!
    \centroiddual(\shift, \weight) \!+\! \sum\nolimits_{i}\! \weight_i \pvol_i ,
  \label{eq:joint-dual}
\end{equation}
where $\smash{\centroiddual(\shift, \weight) = \energy\bigl( \dual(\shift,\weight), \combweight(\shift,\weight) \bigr)}$ is the centroid dual \eqref{eq:centroid-dual}, now with the weights promoted from fixed inputs to unknowns.
The constrained problem thus becomes the {unconstrained} maximization of $\jointdual$ over the space of power diagrams, parametrized by both the shifts $\shift$ and weights $\weight$.
% now with the sites $\qq$ and the prescribed centroids $\tcentroid$ and volumes $\pvol$ as fixed inputs.

Following Appendix~\ref{sec:appendix-joint}, the gradient of $\jointdual$ stacks the centroid and volume residuals
\begin{subequations}
  \label{eq:joint-gradient}
  \begin{empheq}[left={\empheqlbrace\,}]{align}
    \partial_{\shift_i} \jointdual &= 2\,\cellvol_i\,(\tcentroid_i - \centroid_i) , \label{eq:joint-gradient-shift} 
    \\[2mm]
    \partial_{\weight_i} \jointdual &= \pvol_i - \cellvol_i , \label{eq:joint-gradient-weight}
  \end{empheq}
\end{subequations}
so a stationary configuration is one whose cells realize both the prescribed centroids and volumes.
% 
The Hessian of $\jointdual$ is again a facet-based Laplacian, now assembled from the two multiplier types per site.
Its \emph{diagonal} blocks are exactly the two Hessians we have already seen: the shift--shift block reproduces \eqref{eq:centroid-hessian} from \S\ref{sec:centroid-constraints}, and the weight--weight block reproduces the volume Hessian \eqref{eq:newton-step} from \S\ref{sec:background} for the power diagram of the shifted sites, \ie with $\smash{\norm{\qq_i - \qq_j}}$ replaced with $\smash{\norm{\dual_i - \dual_j}}$.
The only new term is the mixed shift--weight block, which takes the form of:
\begin{subequations}
  \label{eq:joint-hessian}
  \begin{empheq}[left={\empheqlbrace\,}]{align}
    \partial^2_{\shift_i \weight_j} \jointdual
      &= \frac{1}{\norm{\dual_i - \dual_j}} \int_{\face_{ij}} (\xx - \tcentroid_i) \, dA,
    \label{eq:joint-hessian-off}
    \\[2mm]
    \partial^2_{\shift_i \weight_i} \jointdual
      &= -\!\sum\nolimits_{j} \partial^2_{\shift_i \weight_j} \jointdual .
    \label{eq:joint-hessian-diag}
  \end{empheq}
\end{subequations}
Since the weight gradient in \eqref{eq:joint-gradient-weight} is the volume residual, the mixed block in \eqref{eq:joint-hessian} is the negated gradient of the cell volumes with respect to the shifts, \ie $\smash{\partial^2_{\shift_i \weight_j} \jointdual = -\partial \cellvol_j / \partial \shift_i}$.
This mixed term can thus be interpreted as a discrete divergence operator as in \cite{deGoes2015PowerParticles}, which quantifies the infinitesimal volume change of the cells as the shifted sites $\dual$ are displaced by infinitesimal changes of the shifts.
Accordingly, its diagonal \eqref{eq:joint-hessian-diag} states that these volume changes sum to zero, since the cells always tile $\domain$.

Unlike $\centroiddual$, the Hessian of $\jointdual$ always has a nontrivial null-space.
Since the weights are now free, each multiplier $\smash{(\shift_i, \weight_i)}$ moves its lifted site freely in $\smash{\R^{d+1}}$, no longer pivoting about $\smash{(\tcentroid_i, \gamma_i)}$.
Every infinitesimal homothety of the shifted weighted sites is thus realized by some perturbation of the multipliers, whether or not the pivots are coplanar.
Solving for $\smash{(\delta\shift_i, \delta\weight_i)}$, as detailed in Appendix~\ref{sec:appendix-joint}, the Hessian annihilates every perturbation of the form
\begin{equation}
  \delta\shift_i = \mu\,\dual_i + \translation , \qquad \delta\weight_i = \mu\,\gamma_i + 2\,\translation\tr \tcentroid_i + c ,
  \label{eq:joint-nullspace}
\end{equation}
for any $\smash{(\mu, \translation, c) \in \R \times \R^d \times \R}$, with $\gamma_i$ as in \eqref{eq:pivot}.
Since the sites $\qq_i$ are fixed and only the shifts and weights vary, this $(d\!+\!2)$-dimensional family is precisely the infinitesimal version of the redundancy of the triplets $\smash{(\qq_i, \shift_i, \weight_i)}$ characterized by \citet[Lemma~1]{Memari2012PrimalDual}.
Its weight update vanishes exactly under \eqref{eq:centroid-kernel}, recovering the Hessian null-space of $\centroiddual$, and for $\mu \!=\! 0$ and $\translation \!=\! 0$ it reduces to the additive constant spanning the Hessian null-space of \eqref{eq:newton-step}.
Therefore, the Hessian of $\jointdual$ is negative definite away from the modes in \eqref{eq:joint-nullspace}.
As in \S\ref{sec:centroid-constraints}, this non-uniqueness concerns the multipliers, not the diagram.
Since every mode in \eqref{eq:joint-nullspace} leaves the cells unchanged, the maximizing multipliers $(\shift, \weight)$, when they exist, form a $(d\!+\!2)$-dimensional family that generates a \emph{single} power diagram.
In terms of the achievable centroids, the dilation $\mu$ and translation $\translation$ move the generators $\dual$ within the normal cone of $\achievable$ at the centroids (\S\ref{sec:background}), while the constant $c$ only affects the weights.

\subsection{Softened Centroid Constraints}
\label{sec:regularized-joint}

Although the jointly constrained problem in \S\ref{sec:joint-constraints} has a unique solution whenever one exists, not every choice of volume and centroid targets is realized by a power diagram.
% since realizable targets must be exposed points of the achievable centroids $\achievable$ 
% 
However, the two constraint types play asymmetric roles in this limitation.
Any positive volumes verifying \eqref{eq:partition-feasibility-vol} are realized by a power diagram for any distinct sites \cite{Aurenhammer1998Minkowski}, so infeasibility arises only once the centroids are prescribed as well, as previously characterized by the set $\achievable$ of achievable centroids in \S\ref{sec:background}.
This asymmetry motivates us to consider partitions of $\domain$ in which the volume constraints remain exact and only the centroid constraints are {softened}.

The most direct way to soften the centroid constraints is to penalize the least-squares centroid residual, as done by \citet{Bourne2025Inverting}, who minimized this penalty over the sites while enforcing the volumes through the weights.
However, as acknowledged by the authors, this objective is non-convex in the sites and its optimization is challenging, with local solutions found by a convex surrogate only under implicit assumptions on the target centroids.
In contrast, we show next that the formulation of \S\ref{sec:joint-constraints} admits a small modification that incorporates the softening while keeping the problem an unconstrained concave maximization.

We start from the problem underlying \S\ref{sec:joint-constraints}, which minimizes the transport cost from sites $\qq$ to cells $\cell$ subject to centroid and volume constraints.
We relax each centroid constraint by a slack variable $\smash{\slack_i \!\in\! \R^d}$ that absorbs its residual, and penalize the squared norm of these slacks in the objective, yielding a modified constrained optimization of the form
\begin{subequations}
  \label{eq:regularized-primal}
  \begin{alignat}{3}
    & \min_{\cell, \slack} \quad & & \sum\nolimits_{i} \int_{\cell_i} \norm{\xx - \qq_i}^2 \, d\xx + \alpha \sum\nolimits_i \pvol_i \norm{\slack_i}^2 \label{eq:regularized-primal-obj} \\
    & \text{s.t.} \quad & \constraintbrace{16pt}{8pt} & \int_{\cell_i} (\xx - \tcentroid_i) \, d\xx = \pvol_i \, \slack_i \ \ \forall i , \label{eq:regularized-primal-cons} \\
    & & & \cellvol_i = \pvol_i \ \ \forall i , \label{eq:regularized-primal-vol}
  \end{alignat}
\end{subequations}
where $\alpha \!>\! 0$ acts as the stiffness of the penalty.
% As special cases, note that $\alpha \!=\! 0$ leaves the slacks free thus reducing to the volume-constrained problem \eqref{eq:volume-constrained}, while $\alpha \!\to\! \infty$ forces $\slack \!=\! 0$ and recovers the joint problem of \S\ref{sec:joint-constraints}.
Introducing the same multipliers as in \S\ref{sec:joint-constraints}, the optimal slacks take the closed form $\smash{\slack_i = -\shift_i / \alpha}$, and substituting them back into the Lagrangian leads to the \emph{softened dual objective}
\begin{equation}
  % \jointdual_{\alpha}(\shift, \weight; \qq, \tcentroid) 
  \jointdual_{\alpha}(\shift, \weight) 
  = 
  % \jointdual(\shift, \weight; \qq, \tcentroid) 
  \jointdual(\shift, \weight) 
  - 
  \frac{1}{\alpha} \sum\nolimits_i \pvol_i \norm{\shift_i}^2 ,
  \label{eq:regularized-dual}
\end{equation}
which only adds a concave quadratic term on the shifts to the joint dual \eqref{eq:joint-dual}.
The softened problem is thus again the {unconstrained} maximization of a concave function over the shifts and weights.

Compared to \S\ref{sec:joint-constraints}, the penalty only modifies the shift derivatives of $\jointdual$.
The weight gradient \eqref{eq:joint-gradient-weight} is unchanged, so the volumes remain exact at a maximizer, while the shift gradient \eqref{eq:joint-gradient-shift} becomes
\begin{equation}
  \partial_{\shift_i} \jointdual_{\alpha} = 2\,\cellvol_i\,(\tcentroid_i - \centroid_i) - \frac{2}{\alpha}\,\pvol_i\,\shift_i .
  \label{eq:regularized-gradient-shift}
\end{equation}
Likewise, the Hessian of \eqref{eq:regularized-dual} keeps the weight--weight and mixed blocks from \eqref{eq:joint-hessian}, while each diagonal shift--shift block is shifted by $\smash{-\tfrac{2}{\alpha}\,\pvol_i\,\Id}$.
% Since the Hessian of $\jointdual$ is negative semi-definite, 
This extra term makes the Hessian of $\jointdual_\alpha$ strictly negative along any perturbation with a nonzero shift component, which removes the translation and dilation modes of \eqref{eq:joint-nullspace}.
The remaining perturbations change only the weights, along which the Hessian reduces to the volume Laplacian \eqref{eq:newton-step} of the shifted sites, so the null-space consists only of the additive constant on the weights.
Consequently, $\jointdual_{\alpha}$ is concave for any $\alpha \!>\! 0$ and admits a unique maximizer for {any} target centroids.
% 
However, its optimal partition is a power diagram only if the generators $\smash{\dual}$ remain pairwise distinct.
Setting \eqref{eq:regularized-gradient-shift} to zero, the optimal shifts are $\smash{\shift_i = \alpha\,(\tcentroid_i - \centroid_i)}$, so they satisfy $\smash{\norm{\shift_i} \!\le\! \alpha \operatorname{diam}(\domain)}$ if the targets $\tcentroid$ are within the convex hull of $\domain$.
By the triangle inequality, the generators then verify 
$\smash{
\norm{\dual_i \!-\! \dual_j} 
% \ge \norm{\qq_i - \qq_j} - \norm{\shift_i} - \norm{\shift_j} 
\ge \norm{\qq_i \!-\! \qq_j} \!-\! 2\alpha \operatorname{diam}(\domain)
}$, so they remain pairwise distinct, and thus the maximizer realizes a power diagram, whenever the closest pair of sites is separated by more than $\smash{2\alpha \operatorname{diam}(\domain)}$, \ie
\begin{equation}
  \alpha < \frac{\min_{i \neq j} \norm{\qq_i - \qq_j}}{2 \operatorname{diam}(\domain)} .
  \label{eq:regularized-distinct}
\end{equation}

At the maximizer, \eqref{eq:regularized-primal-cons} forces $\smash{\slack_i = \centroid_i \!-\! \tcentroid_i}$, and the optimal shifts follow from \eqref{eq:regularized-gradient-shift} as $\smash{\shift_i = \alpha\,(\tcentroid_i \!-\! \centroid_i)}$.
The penalty in \eqref{eq:regularized-primal-obj} is thus the volume-weighted least-squares centroid residual $\smash{\alpha \sum_i \pvol_i \norm{\centroid_i - \tcentroid_i}^2}$, which resembles the residual of \citet{Bourne2025Inverting}, except that they weighted each cell by $\smash{\pvol_i^2}$ instead of $\pvol_i$.
While they minimized their residual through a non-convex search over the sites, we minimize ours through a concave maximization.
% Our softened problem penalizes this same residual combined with the first term of \eqref{eq:regularized-primal-obj} through a concave maximization.
% At the solution, we can further rearrange \eqref{eq:regularized-primal-obj}, up to a constant, as
We can further rearrange the optimal value of the objective \eqref{eq:regularized-primal-obj}, up to a constant, as
\begin{equation}
  \alpha \sum\nolimits_i \pvol_i \norm{\centroid_i - \bigl( \tcentroid_i + \tfrac{1}{\alpha}\,\qq_i \bigr)}^2.
  \label{eq:regularized-projection}
\end{equation}
Therefore, maximizing $\jointdual_\alpha$ returns the partition whose centroids $\centroid$ are the closest to the \emph{pushed targets} $\smash{\tcentroid + \qq/\alpha}$ within $\achievable$ in the volume-weighted least-squares sense.
% , \ie their projection onto $\achievable$.
The generators of the resulting power diagram also follow from \eqref{eq:regularized-gradient-shift} as
\begin{equation}
  \dual_i = \qq_i + \shift_i = \alpha \Bigl( \tcentroid_i + \tfrac{1}{\alpha}\,\qq_i - \centroid_i \Bigr),
  \label{eq:regularized-generators}
\end{equation}
indicating the difference between the pushed targets and their projection, which is an outward normal of $\achievable$ at the centroids $\centroid$ (\S\ref{sec:background}).

% Finally, we note that maximizing $\jointdual_\alpha$ corresponds to a proximal step on $\jointdual$ over the generators $\dual$, centered at the sites $\qq$.
% We can thus recover the maximizer of $\jointdual$ by iterating the maximization of $\jointdual_\alpha$, as in the proximal point algorithm \cite{Rockafellar1976Proximal}.
% Starting from sites at the targets $(\smash{\qq^0 \!=\! \tcentroid})$, we move the sites to the generators after each maximization $(\smash{\qq^{k+1} \!=\! \dual^k})$, and warm-start the next one at zero shifts and the previous shifted weights $(\smash{\weight^{k+1} \!=\! \combweight^k})$, which reproduces the previous power diagram.
% For realizable targets, the sites converge to a positive dilation of the generators realizing $\tcentroid$.
% In particular, a single maximization of $\jointdual_\alpha$ suffices when the sites already generate the power diagram realizing $\tcentroid$ up to a homothety \eqref{eq:lifted-homothety}, since the pushed targets then project onto $\tcentroid$ and \eqref{eq:regularized-generators} returns $\dual = \qq$.
% To keep every iterate a power diagram, we also decrease $\alpha$ at each iteration, if needed, so that it satisfies \eqref{eq:regularized-distinct} for the current sites. % (Lemma~\ref{lem:proximal-halving}).
% We summarize this scheme in Algorithm~\ref{alg:proximal}.

\subsection{Softening Both Constraints}
\label{sec:softened-volume}

The softening of \S\ref{sec:regularized-joint} relaxes the centroid constraints but keeps the volumes exact.
For completeness, we also soften the volume constraints, which yields a single dual objective with one stiffness per constraint type.
% As we show next, this dual is strongly concave in all multipliers and contains every formulation of this section as a limiting case.
To this end, we introduce a scalar slack $\smash{\vslack_i \!\in\! \R}$ per cell that absorbs its volume residual, and penalize these slacks with a stiffness $\beta \!>\! 0$ in the objective \eqref{eq:regularized-primal-obj}, together with a linear term given by reference weights $\smash{\bweight = \{\bweight_i\}_{i=1}^n}$, yielding
\begin{subequations}
  \label{eq:softened-volume-primal}
  \begin{alignat}{3}
    & \min_{\cell, \slack, \vslack} \quad & & 
    \sum\nolimits_{i} \int_{\cell_i} \!\norm{\xx - \qq_i}^2 d\xx \!+\! \pvol_i \bigl( \alpha \norm{\slack_i}^2 + \beta\,\vslack_i^2 - \bweight_i\,\vslack_i \bigr) \label{eq:softened-volume-primal-obj} \\
    & \text{s.t.} \quad & \constraintbrace{16pt}{8pt} & \int_{\cell_i} (\xx - \tcentroid_i) \, d\xx = \pvol_i \, \slack_i \ \ \forall i , \label{eq:softened-volume-primal-cons} \\
    & & & \cellvol_i - \pvol_i = \pvol_i \, \vslack_i \ \ \forall i . \label{eq:softened-volume-primal-vol}
  \end{alignat}
\end{subequations}
% 
Employing multipliers as in \S\ref{sec:joint-constraints}, the optimal volume slacks take the closed form $\smash{\vslack_i \!=\! (\bweight_i - \weight_i) / (2\beta)}$, and substituting them back into the objective leads to the \emph{softened dual objective}
\begin{equation}
  \jointdual_{\alpha\beta}(\shift, \weight) = \jointdual_{\alpha}(\shift, \weight) - \frac{1}{4\beta} \sum\nolimits_i \pvol_i\,(\weight_i - \bweight_i)^2 .
  \label{eq:softened-volume-dual}
\end{equation}
% 
Note that this volume-induced penalty only modifies the weight derivatives of $\jointdual_\alpha$, with \eqref{eq:joint-gradient-weight} becoming
\begin{equation}
  \partial_{\weight_i} \jointdual_{\alpha\beta} = \pvol_i - \cellvol_i - \frac{1}{2\beta}\,\pvol_i\,(\weight_i - \bweight_i) ,
  \label{eq:softened-volume-gradient}
\end{equation}
% 
so the optimal weights follow the equation of state
\begin{equation}
  \weight_i = \bweight_i + 2\beta \Bigl( 1 - \frac{\cellvol_i}{\pvol_i} \Bigr),
  \label{eq:softened-volume-state}
\end{equation}
% 
and act as a pressure with $\bweight$ as a reference pressure, while each diagonal weight--weight block of the Hessian is shifted by $\smash{-\tfrac{1}{2\beta}\,\pvol_i}$.
This removes the additive constant on the weights from the null-space, so $\jointdual_{\alpha\beta}$ is strongly concave and admits a unique maximizer for any target constraints.
Since the volumes are no longer exact, the optimal shifts become $\smash{\shift_i = \alpha\,(\cellvol_i / \pvol_i)(\tcentroid_i - \centroid_i)}$, so the generators $\dual$ remain pairwise distinct if the bound \eqref{eq:regularized-distinct} on $\alpha$ is tightened by the factor $\smash{\min_i \pvol_i / \domainVol}$, while $\beta$ requires no bound, since the weights do not move the shifted sites.

It is worth noting that \eqref{eq:softened-volume-dual} covers the softening of each constraint type alone as limiting cases.
In the limit $\alpha \!\to\! 0$, the centroid slacks become free, which forces the shifts to vanish, and $\jointdual_{\alpha\beta}$ reduces to the softening of the volume dual \eqref{eq:volume-dual} by $\beta$:
\begin{equation}
  \dualenergy_\beta(\weight) = \dualenergy(\weight) - \frac{1}{4\beta} \sum\nolimits_i \pvol_i\,(\weight_i - \bweight_i)^2 .
  \label{eq:softened-volume-only-dual}
\end{equation}
In the limit $\beta \!\to\! 0$, the volume slacks become free instead, which forces the weights to the reference weights $\bweight$, and $\jointdual_{\alpha\beta}$ reduces to the softening of the centroid dual \eqref{eq:centroid-dual} by $\alpha$ for the fixed weights $\bweight$:
\begin{equation}
  \centroiddual_\alpha(\shift) = \centroiddual(\shift) - \frac{1}{\alpha} \sum\nolimits_i \pvol_i \norm{\shift_i}^2 ,
  \label{eq:softened-centroid-dual}
\end{equation}
where $\pvol_i$ only serves as a reference volume weighting the penalty, since the volumes are free.
% 
% The stiffness parameters $(\alpha, \beta)$ thus connect all the formulations of this section: 
Conversely, setting $\beta \!\to\! \infty$ enforces exact volumes and recovers $\jointdual_\alpha$ (\S\ref{sec:regularized-joint}) and $\jointdual$ as $\alpha \!\to\! \infty$ (\S\ref{sec:joint-constraints})
% , while $\alpha \!\to\! 0$ and $\beta \!\to\! 0$ recover the softened duals $\dualenergy_\beta$ and $\centroiddual_\alpha$, respectively, and in turn $\dualenergy$ (\S\ref{sec:background}) and $\centroiddual$ (\S\ref{sec:centroid-constraints}) as their remaining stiffness tends to infinity.
% , which suggests extending the flows of \S\ref{sec:flow} to compressible fluids, which we leave for future work.
% Softening the volumes thus enables an extension of the incompressible flows of \cite{deGoes2015PowerParticles} to the compressible regime, in which the weights play the role of the pressure given by an equation of state (\S\ref{sec:flow-centroids}).

\subsection{Remarks}
\label{sec:remarks}

We close this section with a few remarks on our formulations.
First, we relate our approach to prior work by \citet{Bourne2025Inverting, Bourne2026ConvexOrder}, and then we discuss the feasibility of the prescribed volumes and centroids as well as the existence of solutions.

\paragraph{Link to \citet{Bourne2025Inverting}.}
Our joint problem in \S\ref{sec:joint-constraints} is closely related to the inverse problem considered by \citet{Bourne2025Inverting}, which recovers a power diagram from the volumes and centroids of its cells.
In their setting, the unknowns are the sites of the diagram, so the generators of the sought diagram and the sites used in the objective are the same.
In our formulation, the sites $\qq$ and targets $\tcentroid$ are prescribed independently, hence the generators of our solution are the shifted sites $\dual$ of \eqref{eq:shifted-site}.
% 
The two formulations also handle the volume and centroid constraints differently.
The objective gradient proposed by \citet{Bourne2025Inverting} equals the centroid residual only on the submanifold of weights that enforce the target volumes, so their characterization of solutions holds only once the volume constraints are satisfied, requiring every centroid update to be preceded by a weight solve.
We instead impose both constraints on an equal footing through Lagrange multipliers, the weights for the volumes and the shifts for the centroids.
For every value of the multipliers, the volume and centroid residuals are exactly the gradients of $\jointdual$ in \eqref{eq:joint-gradient}, so our characterization of solutions, as stationary points of a single concave function, holds without first satisfying either constraint.
In particular, the centroid residual can be reduced with the weights held fixed, which is precisely the centroid-only problem of \S\ref{sec:centroid-constraints}.
% 
Yet, the two formulations share the same solutions.
At weights that enforce the volumes, $\jointdual$ reduces to their objective, up to a factor of two and an additive constant. % (Appendix~\ref{sec:appendix-regularized}).
Therefore, our results generalize \citet{Bourne2025Inverting} by recovering their inverse problem when the volumes are enforced, while also admitting arbitrary sites, the centroid-only problem of \S\ref{sec:centroid-constraints}, and the softened extensions of \S\ref{sec:regularized-joint} and \S\ref{sec:softened-volume}.

\paragraph{Link to \citet{Bourne2026ConvexOrder}.}
The progression from \S\ref{sec:joint-constraints} to \S\ref{sec:regularized-joint} also relates to the subsequent work of \citet{Bourne2026ConvexOrder}.
Their method dilates the target centroids $\tcentroid$ by a factor $s \!>\! 1$, projects these dilated targets $s\,\tcentroid$ onto $\achievable$, and takes the residual $s\,\tcentroid \!-\! \centroid$ to their projection $\centroid$ as the generators of a power diagram with centroids $\centroid$.
% This method also applies to unrealizable targets, yielding power diagrams whose centroids approximate them, when the resulting generators are distinct.
% 
Instead, our softening of the centroid constraints in \S\ref{sec:regularized-joint} projects the pushed targets $\smash{\tcentroid \!+\! \qq / \alpha}$ onto $\achievable$. % as in \eqref{eq:regularized-projection}.
In the special case of $\tcentroid_i \!=\! \qq_i$, maximizing $\jointdual_\alpha$ is therefore the same problem as the projection-dilation method of \citet{Bourne2026ConvexOrder} with dilation factor $s \!=\! 1 \!+\! 1/\alpha$, and returns the same power diagram.
While their method is solved by first-order iterations that each requires a full volume solve, our approach offers a second-order Newton's method that updates the shifts and weights simultaneously.
% Moreover, since the sites of $\jointdual_\alpha$ are free, we can employ proximal iterations, as described in \S\ref{sec:regularized-joint}, to recover realizable target centroids exactly, whereas \citet{Bourne2026ConvexOrder} only approaches realizable targets in the limit as $s \!\to\! 1$.

\paragraph{Feasibility.}
Both formulations in \S\ref{sec:centroid-constraints} and \S\ref{sec:joint-constraints} require targets compatible with the conservation identities \eqref{eq:partition-feasibility}, which give simple necessary conditions for feasibility.
% 
In the centroid-only case from \S\ref{sec:centroid-constraints}, the volumes are unknown, and we must check whether positive volumes satisfying \eqref{eq:partition-feasibility} exist, which can be done, \eg, as detailed in Appendix~\ref{sec:appendix-maxent}.
% 
In the joint setting from \S\ref{sec:joint-constraints}, where both volumes and centroids are prescribed, we can simply substitute the data into \eqref{eq:partition-feasibility}.
Note, however, that these tests are only necessary in both settings, \ie a failure proves that the targets are infeasible, but a success does not guarantee that a power diagram realizes them.
In contrast, the softened formulations of \S\ref{sec:regularized-joint} and \S\ref{sec:softened-volume} need no test on the centroids, since the slacks absorb any violation of \eqref{eq:partition-feasibility-centroid}, which leaves only \eqref{eq:partition-feasibility-vol} on the prescribed volumes for $\jointdual_\alpha$, and no condition at all for $\jointdual_{\alpha\beta}$.

\paragraph{Existence.}
Since we restrict the generators to be distinct, our formulations admit a solution only if its energy attains a finite maximum at such a configuration.
We summarize the conditions under which this holds for each formulation. 
% The admissible targets narrow as each formulation prescribes more of the diagram, and widen again once the centroid constraints are softened.
For $\dualenergy$, \citet{Aurenhammer1998Minkowski} showed that, for any distinct generators, there exist weights realizing any positive volumes summing to $\domainVol$, so that no generator is hidden at the maximizer.
For $\centroiddual$, not every prescribed centroids are realizable and the energy may instead approach its supremum as generators merge.
% 
For $\jointdual$, the prescribed centroids must be an exposed point of $\achievable$, which requires, as characterized by \citet{Bourne2025Inverting, Bourne2026ConvexOrder}, that the largest factor $s$ for which the dilated targets $\smash{\domainCentroid + s\,(\tcentroid_i - \domainCentroid)}$ remain in $\achievable$ equals one.
% Therefore, the solution existence for $\jointdual$ adds one scalar condition to \eqref{eq:partition-feasibility} on the spread of the prescribed centroids relative to $\domainCentroid$.
% Otherwise, the energy $\jointdual$ either grows without bound as the shifted sites dilate, for targets outside $\achievable$, or approaches its supremum as they merge, for targets inside $\achievable$.
Finally, $\jointdual_{\alpha}$ trades exact centroids for the volume-weighted projection of the pushed targets $\smash{\tcentroid + \qq/\alpha}$ onto $\achievable$, which always exists and is a power diagram under \eqref{eq:regularized-distinct}.

\end{document}
